

%In this line of my existing work, I developed methods that incorporate domain knowledge to solve scientific problems that have high noise or ill-posed with the data being used, with applications ranging from systems biology, microbiology, to alignment of LLMs. 

%In this line of work, I developed probabilistic machine learning methods that incorporate domain knowledge to solve scientific problems where data are scarce, noisy, or the problems are ill-posed. These methods embed structural priors strongly enough to make inference tractable, yet flexibly enough to handle imperfect or incomplete knowledge. I have applied this approach across diverse domains, including systems biology, microbial biology, and the alignment of LLMs.

Many scientific inference problems are fundamentally ill-posed: multiple underlying mechanisms can explain the same limited or noisy observations. My work develops probabilistic methods that restore identifiability by embedding structured domain knowledge directly into the inference process. Rather than treating domain assumptions as fixed or complete, I model them as soft constraints---structured enough to regularize inference, yet flexible enough to tolerate model misspecification. I have applied this idea in genomics at different levels, including single cell sequencing and large population genetics surveys.
